% Section 4.1, from lectures/L04/appendix/a_kernel_objects.md.

\section{Kernel Objects: Synchronization, Data and Memory}
\label{c4:sec:kernelobjectssynch}\label{c4:app:a}
Chapter~\ref{c3:ch} was about threads on their own: who runs, when, and for how long. This section
is about threads together: how they wait for each other, protect what they share, hand each other
data, and get memory without \code{malloc()}. By the end you should be able to explain every line
the \file{lectures/L04/examples/msgq} and \file{lectures/L04/examples/priority_inversion} examples
print. \secref{c4:app:b} builds on it with two ways of structuring an application.

\subsection{Waiting, and how long for}
Every kernel call that can block takes a \textbf{timeout} as its last argument, of type
\code{k_timeout_t}:

\begin{booktable}{@{}LL@{}}
  \thead{Timeout} & \thead{Meaning} \\
  \midrule
  \code{K_NO_WAIT} & Do not wait. If the call cannot succeed at once, it fails at once. \\
  \code{K_MSEC(n)}, \code{K_SECONDS(n)}, … & Wait at most that long, then fail. \\
  \code{K_FOREVER} & Wait for as long as it takes. \\
\end{booktable}

A call that fails tells you why through its return value, and the two failures mean different
things. \code{-EBUSY} (or \code{-ENOMSG}, or \code{-ENOMEM}, depending on the object) says \emph{it
was not possible right now, and you asked not to wait}. \code{-EAGAIN} says \emph{you waited as long
as you allowed, and it still did not become possible}. Both are ordinary results, not errors to
ignore: a thread that takes a lock with a timeout and does not check whether it got it is a data
race waiting to happen.

\textbf{An interrupt handler must always pass \code{K_NO_WAIT}}, because an ISR has no thread to put
to sleep. The API documentation marks every function an ISR may call with \emph{isr-ok}, and the
ones that take a timeout add \emph{timeout must be set to \code{K_NO_WAIT} if called from ISR}.

\textbf{A thread waiting on an object waits in a queue ordered by priority.} When the object becomes
available, the highest-priority waiter gets it, and among waiters of equal priority, the one that
has waited longest. A thread that becomes ready this way preempts the current thread if it has a
higher priority, exactly as a thread waking from \code{k_sleep()} would.

\subsection{Semaphores}
A \code{k_sem} is a counter with a limit:

\begin{ccode}
K_SEM_DEFINE(data_ready, 0, 1);   /* Initial count 0, limit 1. */

/* In an interrupt handler, or another thread: */
k_sem_give(&data_ready);

/* In the thread that processes the data: */
if (k_sem_take(&data_ready, K_MSEC(100)) == 0) {
	process();
}
\end{ccode}

\noindent \code{k_sem_take()} decrements the count if it is above zero, and otherwise waits. It
returns \code{0} when it got the semaphore, \code{-EBUSY} when it did not and was told not to wait,
and \code{-EAGAIN} when its timeout expired. \code{k_sem_give()} wakes the first waiter if there is
one, and otherwise increments the count, \textbf{but never above the limit}: three gives to a
semaphore with limit 1 and no waiters leave its count at 1. It never blocks, and may be called from
an ISR.

Semaphores are the tool for \textbf{signalling and counting}: \emph{data is ready}, \emph{an event
happened}, \emph{there are N free buffers}. A semaphore with limit 1 signals that something happened
at least once since the last take, not how many times.

\textbf{A semaphore has no owner.} Any thread may give it, whether or not that thread took it, and
the kernel does not know or care which thread ``holds'' it. That is exactly what makes a semaphore
right for signalling, and wrong as a lock, for a reason the next two sections are about.

\subsection{Mutexes}
A \code{k_mutex} is a lock with an owner:

\begin{ccode}
K_MUTEX_DEFINE(bus_lock);

k_mutex_lock(&bus_lock, K_FOREVER);
/* ... use the bus: the critical section ... */
k_mutex_unlock(&bus_lock);
\end{ccode}

\noindent Three things follow from having an owner:

\begin{itemize}
  \item \textbf{Only the owner may unlock it.} \code{k_mutex_unlock()} from another thread returns
    \code{-EPERM}, and on a mutex that is not locked at all, \code{-EINVAL}.
  \item \textbf{It is recursive.} The owner may lock it again without blocking: the kernel counts the locks,
    and the mutex is released when the owner has unlocked it as many times as it locked it. A helper
    that locks the bus can therefore be called from code that already holds the bus lock.
  \item \textbf{It cannot be used in an ISR.} An interrupt handler is not a thread, so it cannot own
    anything. An ISR that needs to share data with a thread uses a spinlock, \code{irq_lock()}, or
    better, hands the data over through a message queue.
\end{itemize}

And one thing that matters more than all three: \textbf{priority inheritance}.

\subsubsection{Priority inversion}
Three threads share one CPU: \code{high}, \code{medium} and \code{low}. \code{low} locks a resource
that \code{high} also needs. Then:

\begin{enumerate}
  \item \code{low} takes the lock and starts its critical section.
  \item \code{high} becomes ready, preempts \code{low}, tries to take the lock, and blocks.
  \item \code{medium} becomes ready. It needs no lock at all. It has a higher priority than \code{low}, so
    it preempts \code{low}, and runs for as long as it has work.
  \item Only when \code{medium} is done does \code{low} get the CPU back, finish its critical section, and
    release the lock. Only then can \code{high} run.
\end{enumerate}

\code{high}, the most important thread in the system, has waited for \code{medium}, a less important
thread it has nothing to do with. That is \textbf{priority inversion}, and what makes it dangerous
is step 3: \code{high}’s wait is no longer bounded by \code{low}’s critical section, which the
designer knows and keeps short, but by the run time of every thread with a priority between the two,
which nobody bounded. The Mars Pathfinder lander reset itself repeatedly in 1997 for exactly this
reason.

The \file{lectures/L04/examples/priority_inversion} example produces it on purpose. With a semaphore
as the lock, \code{high waited 579 ms}: 69 ms of \code{low}’s remaining work plus all 500 ms of
\code{medium}’s.

\subsubsection{Priority inheritance}
When a thread blocks on a locked \code{k_mutex}, the kernel raises the owner’s priority to the
waiter’s, if the waiter’s is higher. The owner now runs at \code{high}’s priority, so \code{medium}
can no longer preempt it; it finishes its critical section and unlocks, and at the unlock the kernel
restores the priority it had when it locked the mutex. With a mutex in the example, \code{high
waited 79 ms}, just the rest of \code{low}’s critical section, and the trace shows \code{low
(priority 5) gives the mutex back}: \code{low} was running at \code{high}’s priority.

\bookfigure{priority_inversion}{The two runs of the \code{priority_inversion} example on one
scale. A solid bar is a thread with the CPU; a dashed one is a thread that is ready or blocked. With
the semaphore, \code{high} waits for work that has nothing to do with the lock; with the mutex, it
waits only for \code{low}'s critical section.}{c4:fig:inversion}

If the waiting thread gives up because its timeout expires, the kernel lowers the owner’s priority
again, to that of the highest-priority thread still waiting, or back to its own.

What priority inheritance does \textbf{not} do:

\begin{itemize}
  \item \textbf{It does not remove the wait.} \code{high} still waits for \code{low}’s critical section.
    Inheritance bounds the wait by the critical section’s length; keeping that short is still the
    designer’s job.
  \item \textbf{It does not cope well with several mutexes.} The owner’s original priority is saved in each
    mutex when it is locked. If a thread holds several mutexes and releases them in a different order
    than its priority was raised, its priority can be restored wrongly. Zephyr’s documentation
    recommends that a thread lock only one mutex at a time when several are shared between threads of
    different priorities.
  \item \textbf{It does not prevent deadlock.} Two threads that lock two mutexes in opposite orders still
    wait for each other forever. Always take nested locks in the same order.
\end{itemize}

\code{CONFIG_PRIORITY_CEILING} caps how high inheritance may raise a priority. Its default, -128,
allows any.

\subsubsection{On \code{native_sim}: making a thread use CPU time}
Priority inversion needs a thread that \emph{computes}: holds the CPU without blocking. On
\code{native_sim}, code runs in zero simulated time, so a loop of instructions takes no time at all,
and \code{medium} would finish before anything could preempt it. The examples and exercises use
\textbf{\code{k_busy_wait()}} instead, which advances simulated time while keeping the CPU, and
during which interrupts are still handled: a thread that becomes ready meanwhile preempts the busy
thread, as it would on a microcontroller. The examples busy-wait one millisecond at a time, because
a busy wait ends at a fixed deadline even if the thread was preempted in between: one call of
\code{k_busy_wait(100000)} would count time spent preempted as work done.

\subsection{Events}
A \code{k_event} is a set of 32 bits that threads can wait on:

\begin{ccode}
K_EVENT_DEFINE(heater_events);

#define EVENT_START BIT(0)
#define EVENT_FAULT BIT(4)

/* From anywhere, including an ISR: */
k_event_post(&heater_events, EVENT_FAULT);

/* In the thread: wait for any of the events, and clear the ones returned. */
uint32_t events = k_event_wait_safe(&heater_events, EVENT_START | EVENT_FAULT,
				   false, K_FOREVER);
\end{ccode}

\noindent \code{k_event_post()} sets bits, \code{k_event_set()} replaces them all, and
\code{k_event_clear()} clears some. \code{k_event_wait()} returns as soon as any of the requested
bits is set, and \code{k_event_wait_all()} when all of them are. Both return the matching bits, or 0
on a timeout. The \code{_safe} versions also clear the bits they return in the same step, so that an
event posted between the wait returning and a separate \code{k_event_clear()} cannot be lost. Events
need \code{CONFIG_EVENTS=y}.

\textbf{An event object is a set of bits, not a queue.} Posting a bit that is already set does
nothing, so three posts of \code{EVENT_START} before the thread runs are one \code{EVENT_START}, and
the order in which different bits were posted is lost. The \file{lectures/L04/examples/smf} example
shows it happening. Use events for \emph{which conditions are true now}, and a message queue when
the count or the order matters.

\subsection{Message queues}
A \code{k_msgq} passes fixed-size messages between threads, by copying them:

\begin{ccode}
struct sample {
	uint32_t seq;
	int64_t timestamp_ms;
};

/* Room for four samples, allocated statically.
 * The last argument is the alignment. */
K_MSGQ_DEFINE(samples, sizeof(struct sample), 4, alignof(struct sample));

/* The producer: */
struct sample s = {.seq = 1, .timestamp_ms = k_uptime_get()};
int err = k_msgq_put(&samples, &s, K_NO_WAIT);

/* The consumer: */
struct sample s;
k_msgq_get(&samples, &s, K_FOREVER);
\end{ccode}

\noindent \code{k_msgq_put()} copies the message into the queue, and \code{k_msgq_get()} copies it
out, first in, first out. Both return \code{0} on success, \code{-ENOMSG} when they could not and
were told not to wait, and \code{-EAGAIN} when their timeout expired. Put may be called from an ISR
with \code{K_NO_WAIT}, which makes a message queue the standard way to hand data from an interrupt
handler to a thread.

Because messages are copied, \textbf{the sender’s variable can go out of scope as soon as the put
returns}, and the queue’s memory is the only memory involved: \code{max_msgs × msg_size} bytes,
allocated by \code{K_MSGQ_DEFINE()} at link time. Keep messages small, a few tens of bytes. Larger
data travels better as a pointer in a small message, or through a FIFO.

When a thread is already waiting in \code{k_msgq_get()}, a put copies the message straight into that
thread and makes it ready, and if it has the higher priority, it runs before the put even returns.
That is why the \file{lectures/L04/examples/msgq} example prints \code{consumed \#1} before
\code{produced \#1}, and a queue that is \code{0/4} right after the first put.

Other calls: \code{k_msgq_peek()} reads the next message without removing it,
\code{k_msgq_put_front()} puts a message at the head, \code{k_msgq_purge()} discards everything and
wakes blocked senders with \code{-ENOMSG}, and \code{k_msgq_num_used_get()} and
\code{k_msgq_num_free_get()} report how full the queue is.

\subsubsection{When the queue is full}
A full queue is a design decision, not an accident, and the producer’s timeout is where it is made:

\begin{itemize}
  \item \textbf{\code{K_FOREVER}: slow the producer down.} Nothing is lost, and the producer runs at the
    consumer’s pace. Right when every message matters and the producer can wait, such as commands.
  \item \textbf{\code{K_NO_WAIT}: drop the new message.} The producer keeps its rate, and must decide what
    to do with the failure: count it, report it. The only option in an ISR.
  \item \textbf{Make room: drop the oldest.} With one producer, \code{k_msgq_get()} with \code{K_NO_WAIT}
    and a second put keep the newest data, which is usually what sensor data wants.
  \item \textbf{A timeout: wait a little.} Absorbs short bursts without letting a stuck consumer stop the
    producer.
\end{itemize}

The example’s two phases show the first two: with \code{K_NO_WAIT} samples \#8, \#10 and \#12 are
dropped, and with \code{K_FOREVER} nothing is, but from \#21 on every put blocks for 150 ms.

\subsection{FIFOs and LIFOs}
A \code{k_fifo} passes data items \textbf{by reference}:

\begin{ccode}
struct packet {
	void *reserved;      /* The first word belongs to the kernel. */
	size_t length;
	uint8_t data[128];
};

K_FIFO_DEFINE(rx_packets);

/* Put never blocks and may be called from an ISR; get returns NULL on timeout. */
k_fifo_put(&rx_packets, packet);
struct packet *p = k_fifo_get(&rx_packets, K_FOREVER);
\end{ccode}

\noindent Nothing is copied: the FIFO links the items together through their \textbf{first word},
which the kernel overwrites, so every item type starts with a pointer-sized member the code never
touches. The item must stay valid until it is taken out, so it cannot be a local variable of the
sender; it comes from a static pool, a memory slab or a heap. A FIFO has no size limit, because the
items bring their own memory, and it never blocks on put. A \code{k_lifo} is the same, last in,
first out.

Use a message queue for small values and a FIFO for large buffers that already live somewhere.

\subsection{Waiting for several things: \code{k_poll()}}
A thread that must react to a command \emph{or} new data cannot block in \code{k_msgq_get()} on one
of them. \code{k_poll()} waits for any of several conditions:

\begin{ccode}
struct k_poll_event events[] = {
	K_POLL_EVENT_INITIALIZER(K_POLL_TYPE_SEM_AVAILABLE,
				 K_POLL_MODE_NOTIFY_ONLY, &data_ready),
	K_POLL_EVENT_INITIALIZER(K_POLL_TYPE_MSGQ_DATA_AVAILABLE,
				 K_POLL_MODE_NOTIFY_ONLY, &commands),
};

while (true) {
	k_poll(events, ARRAY_SIZE(events), K_FOREVER);

	if (events[0].state == K_POLL_STATE_SEM_AVAILABLE) {
		k_sem_take(&data_ready, K_NO_WAIT);
		process_data();
	}
	if (events[1].state == K_POLL_STATE_MSGQ_DATA_AVAILABLE) {
		k_msgq_get(&commands, &command, K_NO_WAIT);
		handle(&command);
	}
	events[0].state = K_POLL_STATE_NOT_READY;
	events[1].state = K_POLL_STATE_NOT_READY;
}
\end{ccode}

\noindent Three rules, all easy to miss:

\begin{itemize}
  \item \textbf{\code{k_poll()} does not take anything.} It only reports that the semaphore was available:
    the thread still takes it, with \code{K_NO_WAIT}, and must handle failing to, since a thread blocked
    in \code{k_sem_take()} has precedence over a polling one.
  \item \textbf{Check every event’s \code{state}}, since several may be ready at once.
  \item \textbf{Reset every \code{state} to \code{K_POLL_STATE_NOT_READY}} before polling again.
\end{itemize}

Polling needs \code{CONFIG_POLL=y}. Besides semaphores and message queues, it can wait on FIFOs and
on a \code{k_poll_signal}, a lightweight flag raised with \code{k_poll_signal_raise()}.

\subsection{Choosing}
\begin{booktable}{@{}LL@{}}
  \thead{You need to} & \thead{Use} \\
  \midrule
  Protect data shared between threads & \code{k_mutex} \\
  Protect data shared with an ISR & A spinlock or \code{irq_lock()}, held briefly; or pass the data, below \\
  Signal that something happened, or count available resources & \code{k_sem} \\
  Wait for any or all of several conditions & \code{k_event}, or \code{k_poll()} \\
  Pass small values, copied, from a thread or an ISR & \code{k_msgq} \\
  Pass large buffers by reference & \code{k_fifo} \\
  Decouple the parts of an application & zbus, in \secref{c4:app:b} \\
\end{booktable}

\textbf{Never use a semaphore as a lock between threads of different priorities.} It works, until a
thread of medium priority arrives.

\subsection{Memory without malloc}
Every object in this section was defined with a \code{K_..._DEFINE()} macro, at file scope, and so
were the threads and their stacks. That is not a coincidence: \textbf{Zephyr applications allocate
statically}, so that the linker knows all the memory the system will ever use, and the question
\emph{can we run out of memory at run time?} is answered before the firmware ships. For the same
reason, \textbf{Zephyr has no system heap by default}: \code{CONFIG_HEAP_MEM_POOL_SIZE} is 0, and a
call to \code{k_malloc()} does not even link:

\begin{console}
undefined reference to `k_malloc'
\end{console}

\noindent When an application does need memory at run time, Zephyr offers two allocators, each
backed by a static region of its own, so that one part of the system running out cannot starve
another.

\subsubsection{Memory slabs}
A \code{k_mem_slab} hands out fixed-size blocks:

\begin{ccode}
/* Four blocks of 64 bytes, aligned to 8. */
K_MEM_SLAB_DEFINE(packet_slab, 64, 4, 8);

void *block;

if (k_mem_slab_alloc(&packet_slab, &block, K_NO_WAIT) == 0) {
	/* ... */
	k_mem_slab_free(&packet_slab, block);
}
\end{ccode}

\noindent Allocation and freeing take constant time, and since every block has the same size, a slab
cannot fragment: after any sequence of allocations and frees, a free block is always usable. When
the slab is empty, \code{k_mem_slab_alloc()} returns \code{-ENOMEM} without waiting, or waits for a
block to be freed. Allocation may be done in an ISR with \code{K_NO_WAIT}. Slabs are the right tool
for many objects of one type: packets, messages, buffers of a fixed size.

\subsubsection{Heaps}
A \code{k_heap} allocates blocks of any size from its own region:

\begin{ccode}
K_HEAP_DEFINE(frame_heap, 1024);

/* NULL if it could not. */
void *frame = k_heap_alloc(&frame_heap, length, K_MSEC(10));
k_heap_free(&frame_heap, frame);
\end{ccode}

\noindent A heap can fragment, so it can fail even when enough bytes are free in total. Size it with
margin, and use it where sizes genuinely vary. \code{k_malloc()} and \code{k_free()} allocate from
the system heap, which exists only when \code{CONFIG_HEAP_MEM_POOL_SIZE} gives it a size.

The C library’s \code{malloc()} is a separate allocator again. On \code{native_sim} it is the host’s
own \code{malloc()}, which never runs out in practice, so an application that uses it will behave
very differently on a board. That is one more reason to stay with the kernel’s allocators.

\subsection{Summary}
\begin{itemize}
  \item Blocking calls take a timeout. \code{K_NO_WAIT} fails at once, a timeout fails with \code{-EAGAIN},
    and an ISR may only use \code{K_NO_WAIT}.
  \item Waiters are served highest priority first, and a woken thread of higher priority runs at once.
  \item A semaphore is a counter without an owner: use it for signalling and counting.
  \item A mutex has an owner, is recursive, cannot be used in an ISR, and inherits priority. Use it for
    locking, keep critical sections short, and hold one mutex at a time.
  \item Priority inversion lets medium-priority threads delay a high-priority one without bound. Priority
    inheritance bounds the delay by the critical section.
  \item On \code{native_sim}, CPU work is simulated with \code{k_busy_wait()}, which lets higher-priority
    threads preempt.
  \item A \code{k_event} holds conditions as bits: posts merge, and order is lost.
  \item A message queue copies small messages, and a FIFO links large items by reference. What a producer
    does with a full queue is a design decision.
  \item \code{k_poll()} waits for several objects, but takes none of them.
  \item Allocate statically. There is no system heap unless one is configured; memory slabs and heaps give
    run-time allocation from bounded regions.
\end{itemize}
